% =========================================================================
%  HDT 1 - Computacion Paralela y Distribuida
%  Informe formal: medicion de rendimiento del conteo de frecuencia
%  Compilar con:  pdflatex informe.tex  (dos veces, por el indice)
% =========================================================================
\documentclass[11pt,letterpaper]{article}

\usepackage[T1]{fontenc}
\usepackage[utf8]{inputenc}
\usepackage[spanish,es-nodecimaldot]{babel}
\usepackage{lmodern}
\usepackage{microtype}
\usepackage[margin=2.6cm,top=2.8cm,bottom=2.8cm]{geometry}
\usepackage{amsmath}
\usepackage{booktabs}
\usepackage{multirow}
\usepackage{array}
\usepackage{xcolor}
\usepackage{graphicx}
\usepackage{float}
\usepackage{caption}
\usepackage{enumitem}
\usepackage{fancyhdr}
\usepackage{titlesec}
\usepackage{listings}
\usepackage{siunitx}
\usepackage[most]{tcolorbox}
\usepackage{tikz}
\usetikzlibrary{shapes.geometric,arrows.meta,positioning,fit,backgrounds,calc}
\sisetup{group-separator={\,},group-minimum-digits=4,output-decimal-marker={.}}
\usepackage[hidelinks,colorlinks=true,linkcolor=azulunc,urlcolor=azulunc,
            citecolor=azulunc]{hyperref}

% ----------------------------------------------------------------- colores
\definecolor{azulunc}{HTML}{1B3A5C}
\definecolor{azulcl}{HTML}{2E6DA4}
\definecolor{grisln}{HTML}{C8D0DA}
\definecolor{grisbg}{HTML}{F4F6F8}
\definecolor{verdeok}{HTML}{1E7A46}
\definecolor{naranja}{HTML}{B45309}
\definecolor{rojoal}{HTML}{A32020}
\definecolor{codecom}{HTML}{6A8759}
\definecolor{codekey}{HTML}{0B5394}

% ----------------------------------------------------------------- titulos
\titleformat{\section}{\normalfont\Large\bfseries\color{azulunc}}
            {\thesection.}{0.6em}{}
\titleformat{\subsection}{\normalfont\large\bfseries\color{azulcl}}
            {\thesubsection.}{0.6em}{}
\titleformat{\subsubsection}{\normalfont\normalsize\bfseries\color{azulcl}}
            {\thesubsubsection.}{0.6em}{}
\titlespacing*{\section}{0pt}{16pt}{8pt}

% ------------------------------------------------------------ encabezados
\pagestyle{fancy}
\fancyhf{}
\renewcommand{\headrulewidth}{0.4pt}
\renewcommand{\footrulewidth}{0.4pt}
\fancyhead[L]{\footnotesize\color{azulunc}Computación Paralela y Distribuida}
\fancyhead[R]{\footnotesize\color{azulunc}Hoja de Trabajo No.~1}
\fancyfoot[C]{\footnotesize\thepage}

% ---------------------------------------------------------------- listings
\lstdefinestyle{c}{
  language=C,
  basicstyle=\ttfamily\scriptsize,
  keywordstyle=\color{codekey}\bfseries,
  commentstyle=\color{codecom}\itshape,
  stringstyle=\color{naranja},
  numbers=left, numberstyle=\tiny\color{gray}, numbersep=8pt,
  backgroundcolor=\color{grisbg},
  frame=single, rulecolor=\color{grisln},
  breaklines=true, showstringspaces=false,
  xleftmargin=16pt, xrightmargin=4pt,
  framexleftmargin=12pt,
  aboveskip=10pt, belowskip=8pt,
  captionpos=b,
}
\lstdefinestyle{sh}{
  basicstyle=\ttfamily\scriptsize,
  backgroundcolor=\color{grisbg},
  frame=single, rulecolor=\color{grisln},
  breaklines=true, showstringspaces=false,
  xleftmargin=4pt, xrightmargin=4pt,
  aboveskip=10pt, belowskip=8pt,
}
\lstset{style=c}

% ------------------------------------------------------------------ cajas
\newtcolorbox{cajadest}[1][azulcl]{
  enhanced, breakable, sharp corners,
  colback=grisbg, colframe=#1,
  boxrule=0pt, leftrule=3pt,
  left=9pt, right=9pt, top=7pt, bottom=7pt,
  width=0.96\linewidth, center,
  before skip=10pt, after skip=10pt,
}
\newcommand{\destacado}[2][azulcl]{%
  \begin{cajadest}[#1]\small #2\end{cajadest}}

\newcolumntype{R}{>{\raggedleft\arraybackslash}p{2.2cm}}
\captionsetup{font=small,labelfont={bf,color=azulunc}}
\setlist{itemsep=2pt,parsep=0pt,topsep=4pt}

% =========================================================================
\begin{document}

% ------------------------------------------------------------------ PORTADA
\begin{titlepage}
\centering
\vspace*{0.3cm}

{\color{azulunc}\rule{\linewidth}{2pt}}\\[6pt]
{\large\bfseries\color{azulunc} COMPUTACIÓN PARALELA Y DISTRIBUIDA}\\[4pt]
{\color{azulunc}\rule{\linewidth}{0.8pt}}

\vspace{1.5cm}

{\Huge\bfseries\color{azulunc} Medición de Rendimiento\\[10pt]
 en el Conteo de Frecuencia\\[10pt] de Palabras}

\vspace{0.7cm}

{\Large\color{azulcl} Comparación entre una implementación secuencial\\[4pt]
 y una implementación paralela con POSIX threads}

\vspace{0.5cm}

{\large\itshape Análisis de tiempo de ejecución, \emph{speedup} y \emph{efficiency}}

\vspace{1.4cm}

{\color{grisln}\rule{0.55\linewidth}{0.6pt}}

\vspace{0.6cm}

{\large\bfseries Integrantes del grupo}

\vspace{0.5cm}

{\large
\begin{tabular}{c}
Pablo Cabrera \\[5pt]
Eliazar Canastuju \\[5pt]
Diego Ramírez \\[5pt]
André Pivaral \\
\end{tabular}}

\vspace{1.0cm}

{\color{grisln}\rule{0.55\linewidth}{0.6pt}}

\vspace{1.1cm}

{\large\bfseries Hoja de Trabajo No.~1}\\[8pt]
{\large 11 de agosto de 2026}

\vspace{0.6cm}
{\color{azulunc}\rule{\linewidth}{2pt}}
\end{titlepage}

% ------------------------------------------------------------------ INDICE
\tableofcontents
\thispagestyle{fancy}
\newpage

% =========================================================== 1. INTRODUCCION
\section{Introducción y objetivos}

El presente informe documenta la medición y comparación del rendimiento de dos
implementaciones del algoritmo de conteo de frecuencia de palabras: una
\textbf{secuencial} y una \textbf{paralela}. Ambas parten del programa desarrollado
en la actividad Corto~02 y fueron instrumentadas para medir el tiempo empleado
específicamente durante la fase de conteo.

Los objetivos concretos del trabajo son:

\begin{enumerate}[label=\textbf{\arabic*.}]
  \item Instrumentar ambas versiones para medir el tiempo de la fase de conteo,
        asegurando que la medición cubra exactamente la misma porción del
        algoritmo en las dos.
  \item Ejecutar cinco corridas por configuración (secuencial, 2~hilos y
        4~hilos) sobre el mismo archivo de texto y obtener el tiempo promedio.
  \item Calcular el \emph{speedup} y la \emph{efficiency} de cada configuración
        paralela.
  \item Analizar el comportamiento del rendimiento al variar la cantidad de
        hilos e identificar las fuentes de \emph{overhead}.
\end{enumerate}

\destacado[verdeok]{\textbf{Criterio metodológico fundamental.} Un valor de
\emph{speedup} solo tiene significado si ambas versiones calculan exactamente lo
mismo y si el cronómetro cubre el mismo trabajo en las dos. Por esta razón, antes
de reportar cualquier medición se verificó la equivalencia de las salidas
(sección~\ref{sec:correctitud}) y se delimitó con precisión la región medida
(sección~\ref{sec:region}).}

% ============================================================== 2. ENTORNO
\section{Entorno experimental}

Todas las mediciones se realizaron en el mismo equipo y bajo las mismas
condiciones, sin otras aplicaciones de carga activa.

\begin{table}[H]
\centering
\caption{Características del entorno de pruebas}
\small
\begin{tabular}{@{}p{5.1cm}p{9.2cm}@{}}
\toprule
\textbf{Elemento} & \textbf{Valor} \\
\midrule
Procesador & Apple M5 (arquitectura ARM64) \\
Núcleos & 10 lógicos: \textbf{4 de rendimiento (P) + 6 de eficiencia (E)} \\
Sistema operativo & macOS 26.5.2 \\
Compilador & Apple clang 21.0.0 \\
Banderas de compilación & \texttt{-O2 -Wall} (secuencial)\newline
                          \texttt{-O2 -Wall -pthread} (paralela) \\
Modelo de paralelismo & POSIX threads (\texttt{pthreads}), memoria compartida \\
Reloj utilizado & \texttt{clock\_gettime(CLOCK\_MONOTONIC)}, resolución de
                  nanosegundos \\
\midrule
Archivo principal & \texttt{archivo.txt}: \num{24530} bytes, \num{4589}
                    palabras, \textbf{904 palabras distintas} \\
Archivo complementario & \texttt{archivo\_grande.txt}: \num{981240} bytes,
                         \num{183560} palabras, 904 distintas \\
\bottomrule
\end{tabular}
\end{table}

\noindent
La arquitectura heterogénea del procesador es relevante para la interpretación de
los resultados: los núcleos P y E \textbf{no tienen el mismo desempeño}, de modo
que el planificador del sistema operativo puede asignar hilos a núcleos de
distinta velocidad.

\subsection{Justificación del archivo complementario}

El enunciado exige emplear el mismo archivo de texto en todas las pruebas,
condición que se cumple dentro de cada experimento. Se incorporó adicionalmente un
segundo experimento sobre un archivo mayor —construido repitiendo cuarenta veces
el mismo texto— por dos motivos:

\begin{itemize}
  \item Con \num{24530} bytes los tiempos son del orden de milisegundos y la
        variabilidad del sistema operativo llega a duplicar la medición entre
        corridas.
  \item Comparar ambos tamaños permite \textbf{aislar el efecto del tamaño del
        problema} sobre el \emph{overhead}, sin modificar una sola línea de
        código. Como se verá en la sección~\ref{sec:analisis}, este contraste
        resultó ser el hallazgo más informativo del trabajo.
\end{itemize}

% =========================================================== 3. DISENO
\section{Diseño de las implementaciones}

\subsection{Estructura de datos común}

Las frecuencias se almacenan en un arreglo dinámico de pares
\emph{(palabra, frecuencia)} que duplica su capacidad al llenarse:

\begin{lstlisting}[caption={Estructura de la tabla de frecuencias}]
typedef struct {
    char   *palabra;      /* copia propia de la palabra */
    size_t  frecuencia;   /* numero de apariciones      */
} EntradaFrecuencia;

typedef struct {
    EntradaFrecuencia *entradas;
    size_t cantidad;      /* entradas utilizadas */
    size_t capacidad;     /* entradas disponibles */
} TablaFrecuencias;
\end{lstlisting}

La inserción recorre linealmente la tabla comparando con \texttt{strcmp}, por lo
que el algoritmo es de complejidad \textbf{$O(n^2)$}: con 904 palabras distintas,
cada inserción puede costar hasta 904 comparaciones. Esta característica,
heredada del programa original, resulta determinante para el experimento, ya que
genera suficiente carga de cómputo como para que la paralelización tenga un
efecto medible.

\subsection{Versión secuencial}

Un único flujo de ejecución recorre el texto carácter por carácter, delimita cada
palabra y la acumula en la tabla:

\begin{lstlisting}[caption={Núcleo del conteo, idéntico en ambas versiones}]
for (i = inicio; i < fin; i++) {
    unsigned char actual = (unsigned char)texto[i];

    if (es_caracter_de_palabra(actual)) {
        palabra[longitud++] = tolower(actual);    /* continua la palabra */
    } else if (longitud > 0) {
        palabra[longitud] = '\0';                 /* termino: registrarla */
        agregar_palabra(tabla, palabra, 1);
        longitud = 0;
    }
}
if (longitud > 0) {          /* si el texto no termina en separador */
    palabra[longitud] = '\0';
    agregar_palabra(tabla, palabra, 1);
}
\end{lstlisting}

\subsubsection{Modificaciones respecto al Corto 02}

La lógica de conteo se conservó sin cambios. Se realizaron tres modificaciones,
todas ellas necesarias para poder medir el tiempo de forma comparable:

\begin{enumerate}[label=\textbf{(\alph*)}]
  \item \textbf{Separación de la E/S respecto del cómputo.} La versión original
        leía con \texttt{fgetc()} carácter por carácter \emph{mientras} contaba.
        Ahora el archivo se carga íntegramente en memoria \emph{antes} de iniciar
        el cronómetro. De no hacerlo, la medición incluiría el acceso a disco en
        lugar de reflejar únicamente el cómputo.
  \item \textbf{Aislamiento del conteo en una función parametrizada.}
        \texttt{contar\_palabras(texto, inicio, fin, tabla)} opera sobre un rango.
        La versión secuencial la invoca una vez sobre el texto completo; en la
        paralela, cada hilo la invoca sobre su bloque. Se trata literalmente de
        \textbf{la misma función en ambas versiones}, lo que garantiza que las
        diferencias de tiempo provengan del paralelismo y no de un cambio de
        algoritmo.
  \item \textbf{Parámetro \texttt{veces} en \texttt{agregar\_palabra}.} La versión
        secuencial siempre pasa el valor 1. El parámetro habilita la fusión de
        conteos parciales en la versión paralela, donde es necesario acumular de
        una sola vez varias apariciones.
\end{enumerate}

\subsection{Versión paralela}

\subsubsection{Estrategia de descomposición}

La descomposición es \textbf{por datos}: el texto se divide en $N$ bloques
contiguos y cada hilo procesa uno. La figura~\ref{fig:flujo} presenta el flujo
completo, incluyendo la delimitación de la región medida.

\begin{figure}[H]
\centering
\begin{tikzpicture}[
  node distance=6mm,
  caja/.style={draw,rounded corners=2pt,minimum height=8mm,align=center,
               font=\footnotesize,inner sep=4pt},
  gris/.style={caja,draw=gray!70,fill=gray!8,text width=68mm},
  azul/.style={caja,draw=azulcl,fill=azulcl!8,text width=68mm},
  verde/.style={caja,draw=verdeok,fill=verdeok!8,text width=25mm,
                minimum height=14mm},
  naranja/.style={caja,draw=naranja,fill=naranja!8,text width=68mm},
  fl/.style={-{Latex[length=2mm]},gray!80,thick},
  ln/.style={gray!80,thick}
]

\node[gris] (leer) {1.\ \ Leer el archivo completo a memoria};
\node[right=3mm of leer,font=\tiny\itshape,gray] {fuera de la medición};

\coordinate (c1) at ($(leer.south)+(0,-4mm)$);
\draw[fl] (leer.south) -- (c1);
\draw[verdeok,thick,dashed] ($(c1)+(-78mm,-1.5mm)$) -- ($(c1)+(78mm,-1.5mm)$)
      node[midway,fill=white,inner sep=2pt,font=\tiny\bfseries,text=verdeok]
      {INICIA EL CRONÓMETRO};

\node[azul,below=9mm of leer] (dividir)
     {2.\ \ Dividir el texto en $N$ bloques\\[1pt]
      \tiny (cada corte se desplaza hasta un separador)};

\coordinate (bus) at ($(dividir.south)+(0,-7mm)$);
\draw[ln] (dividir.south) -- (bus);

\node[verde,below=20mm of dividir.south,xshift=-49mm] (h0)
     {\textbf{Hilo 0}\\[2pt]bloque 0\\tabla privada};
\node[verde,right=6mm of h0] (h1) {\textbf{Hilo 1}\\[2pt]bloque 1\\tabla privada};
\node[verde,right=6mm of h1] (h2) {\textbf{Hilo 2}\\[2pt]bloque 2\\tabla privada};
\node[verde,right=6mm of h2] (h3) {\textbf{Hilo $N{-}1$}\\[2pt]bloque $N{-}1$\\tabla privada};

\draw[ln] (h0.north|-bus) -- (h3.north|-bus);
\foreach \n in {h0,h1,h2,h3} { \draw[fl] (\n.north|-bus) -- (\n.north); }

\node[above=1mm of h1.north east,xshift=3mm,font=\scriptsize\bfseries,
      text=azulcl,yshift=4mm,fill=white,inner sep=3pt]
     {3.\ \ Cada hilo cuenta su bloque en su \emph{propia} tabla};

\coordinate (bus2) at ($(h0.south)+(0,-6mm)$);
\foreach \n in {h0,h1,h2,h3} { \draw[ln] (\n.south) -- (\n.south|-bus2); }
\draw[ln] (h0.south|-bus2) -- (h3.south|-bus2);

\node[naranja,below=14mm of h1.south east,xshift=3mm] (join)
     {4.\ \ \texttt{pthread\_join}: esperar a \textbf{todos} los hilos};
\draw[fl] ($(join.north)+(0,6mm)$) -- (join.north);

\node[naranja,below=6mm of join] (fusion)
     {5.\ \ Fusionar las $N$ tablas privadas en una sola\\[1pt]
      \tiny (porción secuencial: el techo de Amdahl)};
\draw[fl] (join.south) -- (fusion.north);

\coordinate (c2) at ($(fusion.south)+(0,-4mm)$);
\draw[rojoal,thick,dashed] ($(c2)+(-78mm,0)$) -- ($(c2)+(78mm,0)$)
      node[midway,fill=white,inner sep=2pt,font=\tiny\bfseries,text=rojoal]
      {TERMINA EL CRONÓMETRO};

\node[gris,below=9mm of fusion] (salida)
     {6.\ \ Ordenar alfabéticamente e imprimir};
\node[right=3mm of salida,font=\tiny\itshape,gray] {fuera de la medición};

\end{tikzpicture}
\caption{Flujo de la versión paralela. Únicamente la región comprendida entre
         las líneas punteadas forma parte de la medición.}
\label{fig:flujo}
\end{figure}

\subsubsection{División del texto sin fragmentar palabras}

Si los cortes se calcularan por simple división aritmética, uno de ellos podría
caer en medio de una palabra, que quedaría partida entre dos hilos y produciría
dos tokens inexistentes. Para evitarlo, cada corte se desplaza hacia adelante
hasta encontrar un separador:

\begin{lstlisting}[caption={Ajuste de los cortes a frontera de palabra}]
corte = tamano * (i + 1) / num_hilos;

/* Avanzar hasta el siguiente separador para no partir palabras. */
while (corte < tamano && es_caracter_de_palabra((unsigned char)texto[corte])) {
    corte++;
}
\end{lstlisting}

Como consecuencia, los bloques no resultan exactamente del mismo tamaño. Este
desbalance residual constituye una de las fuentes de \emph{overhead} identificadas
en la sección~\ref{sec:overhead}.

\subsubsection{Tablas privadas por hilo}

Cada hilo acumula sus resultados en una tabla a la que ningún otro hilo accede:

\begin{lstlisting}[caption={Descriptor de tarea asignado a cada hilo}]
typedef struct {
    const char *texto;
    size_t inicio, fin;       /* su bloque: [inicio, fin) */
    TablaFrecuencias tabla;   /* tabla privada del hilo   */
    int ok;
} TareaHilo;
\end{lstlisting}

Esta es la decisión de diseño más determinante del programa. La alternativa
—una única tabla compartida protegida por un \emph{mutex}— habría sido
contraproducente: dado que prácticamente toda la ejecución consiste en insertar
palabras en la tabla, los hilos permanecerían la mayor parte del tiempo
esperando el candado. El programa quedaría serializado en la práctica y, sumando
el costo de adquirir y liberar el \emph{mutex} millones de veces, resultaría
previsiblemente \textbf{más lento que la versión secuencial}.

\destacado[verdeok]{Al asignar datos privados a cada hilo, \textbf{la fase paralela
no contiene ni un solo \emph{lock}}. Todo el costo de coordinación se concentra en
un único punto: la reducción final. Esto convierte al programa en un caso
particularmente claro para observar la ley de Amdahl.}

\subsubsection{Sincronización y reducción}

\begin{lstlisting}[caption={Sincronización mediante \texttt{join} y reducción de las tablas}]
for (i = 0; i < num_hilos; i++) {
    pthread_join(hilos[i], NULL);      /* esperar a que todos terminen */
}

/* Reduccion: volcar las tablas 1..N-1 sobre la del hilo 0. */
for (i = 1; i < num_hilos; i++) {
    for (j = 0; j < tareas[i].tabla.cantidad; j++) {
        agregar_palabra(&tareas[0].tabla,
                        tareas[i].tabla.entradas[j].palabra,
                        tareas[i].tabla.entradas[j].frecuencia);
    }
}
\end{lstlisting}

La reducción la ejecuta \textbf{un solo hilo} mientras los restantes ya han
finalizado y permanecen ociosos. Constituye la porción estrictamente secuencial
del algoritmo y, adicionalmente, \textbf{su costo crece con $N$}: con 2 hilos debe
fusionarse una tabla; con 4 hilos, tres.

% ====================================================== 4. METODOLOGIA
\section{Metodología de medición}

\subsection{Delimitación de la región medida}
\label{sec:region}

Se empleó \texttt{clock\_gettime} con el reloj \texttt{CLOCK\_MONOTONIC}, que
ofrece resolución de nanosegundos y, a diferencia del reloj de pared, no
retrocede ante ajustes de hora del sistema.

\begin{table}[H]
\centering
\caption{Delimitación de la región medida en ambas versiones}
\small
\begin{tabular}{@{}p{3.2cm}p{5.3cm}p{5.8cm}@{}}
\toprule
 & \textbf{Versión secuencial} & \textbf{Versión paralela} \\
\midrule
\textbf{Fuera} (antes) & Lectura del archivo a memoria &
  Lectura del archivo a memoria \\
\addlinespace
\textbf{Dentro} & Recorrido del texto y llenado de la tabla &
  División del texto \newline
  Creación de los hilos \newline
  Conteo parcial concurrente \newline
  \texttt{pthread\_join} \newline
  \textbf{Reducción de las tablas} \\
\addlinespace
\textbf{Fuera} (después) & Ordenamiento e impresión &
  Ordenamiento e impresión \\
\bottomrule
\end{tabular}
\end{table}

\destacado{\textbf{La totalidad del \emph{overhead} del paralelismo se encuentra
dentro de la medición}: el reparto del texto, la creación de los hilos, la espera
en la barrera y la reducción final. No se midió únicamente la fase favorable.}

\subsection{Verificación de correctitud}
\label{sec:correctitud}

Se comprobó que la versión paralela produce exactamente la misma tabla de
frecuencias que la secuencial. La verificación está automatizada en el script
\texttt{verificar.sh}, que devuelve código de salida 0 o 1.

\paragraph{Prueba 1: variación en la cantidad de hilos.}
Se comparó la salida completa (904 palabras con sus frecuencias) contra la
referencia secuencial empleando \textbf{1, 2, 3, 4, 5, 7, 8, 16 y 32 hilos}. Se
incluyeron deliberadamente números impares y no divisores del tamaño del archivo,
con el fin de forzar cortes en posiciones irregulares. El resultado fue
\textbf{idéntico en las nueve configuraciones}.

\paragraph{Prueba 2: casos borde.}
Cada caso se contrastó con 1, 2, 4 y 8 hilos.

\begin{table}[H]
\centering
\caption{Casos borde verificados}
\small
\begin{tabular}{@{}p{4.4cm}R p{5.6cm}c@{}}
\toprule
\textbf{Caso} & \textbf{Bytes} & \textbf{Aspecto que ejercita} &
\textbf{Resultado} \\
\midrule
Archivo vacío & 0 & Todos los bloques vacíos & Idéntico \\
Una sola palabra & 4 & Menos palabras que hilos & Idéntico \\
Solo espacios y tabuladores & 7 & Ausencia total de palabras & Idéntico \\
Dos palabras de un carácter & 3 & Más hilos que bytes & Idéntico \\
Palabras repetidas & 16 & Acumulación entre bloques & Idéntico \\
Mayúsculas y minúsculas & 21 & Normalización a minúscula & Idéntico \\
Sin salto de línea final & 15 & Palabra al borde del bloque & Idéntico \\
\bottomrule
\end{tabular}
\end{table}

Los casos de 3 y 4 bytes son los más exigentes: con 8 hilos varios reciben
bloques vacíos y el ajuste de frontera puede producir bloques de longitud cero.
El programa los procesa sin fallar y sin perder ni duplicar palabras.

\subsection{Justificación de las cinco ejecuciones}

Los tiempos presentan una dispersión considerable entre corridas. Las cinco
mediciones secuenciales sobre el archivo de 24~KB oscilaron entre
\num{0.002447}~s y \num{0.004826}~s, es decir, \textbf{casi el doble de diferencia
con el mismo código y el mismo archivo}. Las causas son ajenas al programa: otros
procesos compitiendo por el CPU, decisiones del planificador, variación dinámica
de la frecuencia del procesador y efectos de calentamiento de caché. Promediar
cinco corridas mitiga esta variabilidad.

% ====================================================== 5. RESULTADOS
\section{Resultados experimentales}

\subsection{Tiempos de ejecución}

\begin{table}[H]
\centering
\caption{Tiempos de ejecución sobre \texttt{archivo.txt} (\num{24530} bytes)}
\small
\begin{tabular}{@{}cRRR@{}}
\toprule
\textbf{Núm.} & \textbf{T. secuencial (s)} & \textbf{T. paralela\newline 2 hilos (s)}
& \textbf{T. paralela\newline 4 hilos (s)} \\
\midrule
1 & 0.004635 & 0.003043 & 0.002097 \\
2 & 0.004826 & 0.002977 & 0.001848 \\
3 & 0.003230 & 0.002034 & 0.001863 \\
4 & 0.002752 & 0.001510 & 0.001490 \\
5 & 0.002447 & 0.001376 & 0.001426 \\
\midrule
\textbf{Promedio} & \textbf{0.003578} & \textbf{0.002188} & \textbf{0.001745} \\
\bottomrule
\end{tabular}
\end{table}

\begin{table}[H]
\centering
\caption{Tiempos de ejecución sobre \texttt{archivo\_grande.txt} (\num{981240} bytes)}
\small
\begin{tabular}{@{}cRRR@{}}
\toprule
\textbf{Núm.} & \textbf{T. secuencial (s)} & \textbf{T. paralela\newline 2 hilos (s)}
& \textbf{T. paralela\newline 4 hilos (s)} \\
\midrule
1 & 0.078007 & 0.046513 & 0.024097 \\
2 & 0.069338 & 0.047642 & 0.026287 \\
3 & 0.077322 & 0.041301 & 0.022825 \\
4 & 0.065508 & 0.038404 & 0.023097 \\
5 & 0.067785 & 0.037561 & 0.027518 \\
\midrule
\textbf{Promedio} & \textbf{0.071592} & \textbf{0.042284} & \textbf{0.024765} \\
\bottomrule
\end{tabular}
\end{table}

\subsection{Cálculo del \emph{speedup}}

El \emph{speedup} se define como el cociente entre el tiempo de la versión
secuencial y el de la versión paralela:

\begin{equation}
S_p = \frac{T_{\text{secuencial}}}{T_{\text{paralelo}}}
\end{equation}

\begin{table}[H]
\centering
\caption{\emph{Speedup} por corrida y a partir de los tiempos promedio}
\small
\begin{tabular}{@{}cRR@{\hspace{1.4cm}}cRR@{}}
\toprule
\multicolumn{3}{c}{\textbf{\texttt{archivo.txt} (24 KB)}} &
\multicolumn{3}{c}{\textbf{\texttt{archivo\_grande.txt} (958 KB)}} \\
\cmidrule(r){1-3}\cmidrule(l){4-6}
\textbf{Núm.} & \textbf{2 hilos} & \textbf{4 hilos} &
\textbf{Núm.} & \textbf{2 hilos} & \textbf{4 hilos} \\
\midrule
1 & 1.5232 & 2.2103 & 1 & 1.6771 & 3.2372 \\
2 & 1.6211 & 2.6115 & 2 & 1.4554 & 2.6377 \\
3 & 1.5880 & 1.7338 & 3 & 1.8722 & 3.3876 \\
4 & 1.8225 & 1.8470 & 4 & 1.7058 & 2.8362 \\
5 & 1.7783 & 1.7160 & 5 & 1.8047 & 2.4633 \\
\midrule
\textbf{Prom.} & \textbf{1.6353} & \textbf{2.0504} &
\textbf{Prom.} & \textbf{1.6931} & \textbf{2.8909} \\
\bottomrule
\end{tabular}
\end{table}

\noindent Cálculo a partir de los tiempos promedio del archivo principal:

\begin{align}
S_2 &= \frac{0.003578}{0.002188} = \mathbf{1.6353} \\[2pt]
S_4 &= \frac{0.003578}{0.001745} = \mathbf{2.0504}
\end{align}

\subsection{Cálculo de la \emph{efficiency}}

La \emph{efficiency} expresa el aprovechamiento promedio de cada hilo:

\begin{equation}
E_p = \frac{S_p}{p}
\end{equation}

\begin{table}[H]
\centering
\caption{\emph{Efficiency} por corrida y a partir de los tiempos promedio}
\small
\begin{tabular}{@{}cRR@{\hspace{1.4cm}}cRR@{}}
\toprule
\multicolumn{3}{c}{\textbf{\texttt{archivo.txt} (24 KB)}} &
\multicolumn{3}{c}{\textbf{\texttt{archivo\_grande.txt} (958 KB)}} \\
\cmidrule(r){1-3}\cmidrule(l){4-6}
\textbf{Núm.} & \textbf{2 hilos} & \textbf{4 hilos} &
\textbf{Núm.} & \textbf{2 hilos} & \textbf{4 hilos} \\
\midrule
1 & 0.7616 & 0.5526 & 1 & 0.8386 & 0.8093 \\
2 & 0.8105 & 0.6529 & 2 & 0.7277 & 0.6594 \\
3 & 0.7940 & 0.4334 & 3 & 0.9361 & 0.8469 \\
4 & 0.9113 & 0.4617 & 4 & 0.8529 & 0.7091 \\
5 & 0.8892 & 0.4290 & 5 & 0.9023 & 0.6158 \\
\midrule
\textbf{Prom.} & \textbf{0.8176} & \textbf{0.5126} &
\textbf{Prom.} & \textbf{0.8466} & \textbf{0.7227} \\
\bottomrule
\end{tabular}
\end{table}

\noindent Cálculo a partir de los tiempos promedio del archivo principal:

\begin{align}
E_2 &= \frac{1.6353}{2} = \mathbf{0.8176} \quad \text{(81.76\,\%)} \\[2pt]
E_4 &= \frac{2.0504}{4} = \mathbf{0.5126} \quad \text{(51.26\,\%)}
\end{align}

\subsection{Tabla final de resultados}

\begin{table}[H]
\centering
\caption{Resumen consolidado de resultados}
\small
\begin{tabular}{@{}llRRR@{}}
\toprule
\textbf{Archivo} & \textbf{Configuración} & \textbf{T. promedio (s)} &
\textbf{Speedup} & \textbf{Efficiency} \\
\midrule
\multirow{3}{*}{\shortstack[l]{\texttt{archivo.txt}\\ (24 KB)}}
  & Secuencial & 0.003578 & 1.0000 & 100.00\,\% \\
  & 2 hilos    & 0.002188 & 1.6353 & \textbf{81.76\,\%} \\
  & 4 hilos    & 0.001745 & \textbf{2.0504} & 51.26\,\% \\
\midrule
\multirow{3}{*}{\shortstack[l]{\texttt{archivo\_grande.txt}\\ (958 KB)}}
  & Secuencial & 0.071592 & 1.0000 & 100.00\,\% \\
  & 2 hilos    & 0.042284 & 1.6931 & \textbf{84.66\,\%} \\
  & 4 hilos    & 0.024765 & \textbf{2.8909} & 72.27\,\% \\
\bottomrule
\end{tabular}
\end{table}

\begin{figure}[H]
\centering
\begin{tikzpicture}[x=1.02cm,y=0.62cm]
  % ejes
  \draw[gray!50] (0,0) -- (0,7.4);
  \foreach \x in {0,1,2,3,4} {
    \draw[gray!25,dashed] (\x,0) -- (\x,7.2);
    \node[font=\scriptsize,gray,below] at (\x,0) {\x};
  }
  \draw[gray!50] (0,0) -- (4.4,0);

  % barras: (fila, valor, color, etiqueta)
  \foreach \i/\v/\c/\lbl in {
    6/2.00/gray!45/{Ideal, 2 hilos},
    5/1.6353/azulcl/{24 KB, 2 hilos},
    4/1.6931/verdeok/{958 KB, 2 hilos},
    2.6/4.00/gray!45/{Ideal, 4 hilos},
    1.6/2.0504/azulcl/{24 KB, 4 hilos},
    0.6/2.8909/verdeok/{958 KB, 4 hilos}}
  {
    \fill[\c] (0,\i-0.28) rectangle (\v,\i+0.28);
    \node[font=\scriptsize,right] at (\v+0.08,\i) {\lbl~(\v$\times$)};
  }
  \node[font=\scriptsize\bfseries,gray,below] at (2.2,-0.55)
       {\emph{Speedup} obtenido frente al ideal};
\end{tikzpicture}
\caption{Comparación del \emph{speedup} obtenido contra el \emph{speedup} lineal
         ideal. La brecha se ensancha al pasar de 2 a 4 hilos, y es
         considerablemente menor con el archivo de mayor tamaño.}
\label{fig:speedup}
\end{figure}

% ====================================================== 6. ANALISIS
\section{Análisis de resultados}
\label{sec:analisis}

\subsection{¿Con qué cantidad de hilos se obtuvo el mayor \emph{speedup}?}

Con \textbf{4 hilos} en ambos archivos. En \texttt{archivo.txt} el \emph{speedup}
pasó de 1.6353 (2 hilos) a 2.0504 (4 hilos), y en \texttt{archivo\_grande.txt} de
1.6931 a 2.8909. El mejor resultado absoluto fue \textbf{$2.8909\times$ con 4
hilos sobre el archivo de mayor tamaño}.

No obstante, el incremento de 2 a 4 hilos rindió considerablemente menos de lo
esperado: duplicar la cantidad de hilos debería duplicar el \emph{speedup}, pero
en el archivo pequeño únicamente lo aumentó un 25\,\% ($1.64 \rightarrow 2.05$) y
en el grande un 71\,\% ($1.69 \rightarrow 2.89$).

\subsection{¿Con qué configuración se obtuvo la mayor \emph{efficiency}?}

Con \textbf{2 hilos}: 81.76\,\% en el archivo pequeño y 84.66\,\% en el grande. En
ambas pruebas la eficiencia con 2 hilos superó a la de 4 hilos. El valor más bajo
correspondió a \textbf{4 hilos sobre el archivo pequeño (51.26\,\%)}, donde
prácticamente la mitad de la capacidad de cómputo asignada se consumió en
\emph{overhead}.

\destacado[naranja]{Se observa aquí una tensión característica del cómputo
paralelo: \textbf{la configuración más rápida (4 hilos) no es la más eficiente
(2 hilos)}. Con 2 hilos se aprovecha mejor cada núcleo asignado; con 4 hilos se
termina antes, pero a costa de recursos ociosos.}

\subsection{¿El \emph{speedup} aumentó proporcionalmente?}

\textbf{No.} El crecimiento fue claramente \textbf{sublineal} y se alejó del ideal
conforme se agregaron hilos.

\begin{table}[H]
\centering
\caption{Comparación con el \emph{speedup} lineal ideal}
\small
\begin{tabular}{@{}cRRR@{}}
\toprule
\textbf{Hilos} & \textbf{Speedup ideal} & \textbf{Obtenido (24 KB)} &
\textbf{Obtenido (958 KB)} \\
\midrule
2 & 2.00 & 1.64 \ (82\,\% del ideal) & 1.69 \ (85\,\% del ideal) \\
4 & 4.00 & 2.05 \ (51\,\% del ideal) & 2.89 \ (72\,\% del ideal) \\
\bottomrule
\end{tabular}
\end{table}

Con 2 hilos se alcanzó entre el 82\,\% y el 85\,\% del ideal, mientras que con 4
hilos únicamente entre el 51\,\% y el 72\,\%. La brecha respecto al
\emph{speedup} lineal \textbf{se ensancha} conforme aumenta el número de hilos,
comportamiento que corresponde exactamente a lo que predice la \textbf{ley de
Amdahl}: la porción secuencial del programa —en este caso la fusión de las tablas
privadas— impone un techo al \emph{speedup} alcanzable con independencia de
cuántos hilos se agreguen.

\subsection{¿La \emph{efficiency} aumentó o disminuyó?}

\textbf{Disminuyó de forma consistente} en ambas pruebas:

\begin{itemize}
  \item \texttt{archivo.txt}: 81.76\,\% $\rightarrow$ 51.26\,\%
        (\textbf{caída de 30.5 puntos}).
  \item \texttt{archivo\_grande.txt}: 84.66\,\% $\rightarrow$ 72.27\,\%
        (\textbf{caída de 12.4 puntos}).
\end{itemize}

La razón es que el \emph{overhead} —creación de hilos, reparto del texto y
reducción final— \textbf{crece con el número de hilos}, mientras que el trabajo
útil asignado a cada hilo \textbf{se reduce}. La relación entre trabajo útil y
\emph{overhead} empeora en cada incremento.

Resulta particularmente revelador comparar ambas caídas: con el archivo grande la
pérdida de eficiencia fue \textbf{menos de la mitad} que con el pequeño, dado que
existe considerablemente más trabajo real entre el cual distribuir el mismo costo
fijo de coordinación.

\subsection{Factores que pudieron afectar el tiempo de ejecución}

\begin{enumerate}[label=\textbf{\arabic*.}]
  \item \textbf{Tamaño del archivo} (factor determinante). Fue el efecto más
        marcado del experimento. Con 24~KB el trabajo asignado a cada hilo es tan
        reducido que el costo de crear hilos y fusionar tablas pesa casi tanto
        como el conteo mismo ($E_4 = 51\,\text{\%}$). Al pasar a 958~KB —cuarenta veces
        más texto— la eficiencia con 4 hilos ascendió al 72\,\% \textbf{sin
        modificar una sola línea de código}.
  \item \textbf{Arquitectura heterogénea del procesador.} El Apple~M5 dispone de
        4 núcleos de rendimiento y 6 de eficiencia, que no ofrecen el mismo
        desempeño. Si el planificador ubica algún hilo en un núcleo E, dicho hilo
        tarda más y los restantes deben \textbf{esperarlo en el
        \texttt{pthread\_join}}. El tiempo total lo determina el hilo más lento,
        no el promedio.
  \item \textbf{Variabilidad del sistema operativo.} Los tiempos de una misma
        configuración varían de forma apreciable, con diferencias de hasta el
        doble entre la corrida más lenta y la más rápida, atribuibles a otros
        procesos, al planificador y a los cambios dinámicos de frecuencia.
  \item \textbf{Efectos de caché y calentamiento.} Las primeras corridas
        resultaron sistemáticamente más lentas (\num{0.004635}~s la primera
        frente a \num{0.002447}~s la quinta).
  \item \textbf{Desbalance de carga.} Los bloques se reparten por bytes y no por
        palabras, de modo que los hilos no reciben exactamente la misma cantidad
        de trabajo.
  \item \textbf{Complejidad $O(n^2)$ del algoritmo de tabla.} Cada inserción
        realiza búsqueda lineal. Este costo cuadrático es, paradójicamente, lo
        que hace que la paralelización resulte provechosa: proporciona suficiente
        cómputo que distribuir.
\end{enumerate}

\subsection{Fuentes de \emph{overhead} identificadas}
\label{sec:overhead}

\begin{table}[H]
\centering
\caption{Fuentes de \emph{overhead}, ordenadas por impacto}
\small
\begin{tabular}{@{}p{4.3cm}p{10.2cm}@{}}
\toprule
\textbf{Fuente} & \textbf{Naturaleza del costo} \\
\midrule
\textbf{Reducción de las tablas} \newline \emph{(la principal)} &
  Es 100\,\% secuencial y crece con $N$: con 4 hilos deben fusionarse 3 tablas de
  hasta 904 palabras, y cada inserción realiza búsqueda lineal $O(904)$.
  Constituye el techo de Amdahl del programa. \\
\addlinespace
\textbf{Creación y destrucción de hilos} &
  Cada \texttt{pthread\_create} cuesta decenas de microsegundos. Sobre un trabajo
  total de 3.5~ms, cuatro hilos representan un costo fijo no despreciable. \\
\addlinespace
\textbf{Barrera implícita del \texttt{join}} &
  El tiempo lo determina el hilo más lento; los restantes quedan ociosos. Se
  agrava con el desbalance de carga y con los núcleos heterogéneos. \\
\addlinespace
\textbf{Trabajo duplicado en memoria} &
  Con 4 hilos coexisten 4 tablas de hasta 904 entradas: cuádruple uso de memoria
  y cuádruples llamadas a \texttt{malloc}/\texttt{realloc}. \\
\addlinespace
\textbf{Contención en el asignador} &
  El asignador del sistema emplea \emph{locks} internos; varios hilos solicitando
  memoria simultáneamente se serializan de forma parcial. \\
\addlinespace
\textbf{Cálculo de los cortes} &
  Costo menor, pero corresponde a trabajo que la versión secuencial no realiza. \\
\addlinespace
\textbf{\emph{False sharing}} &
  Estructuras \texttt{TareaHilo} contiguas pueden compartir línea de caché,
  provocando invalidaciones cruzadas. \\
\bottomrule
\end{tabular}
\end{table}

Conviene destacar que la fuente de \emph{overhead} que se \textbf{evitó} explica
también los buenos resultados con 2 hilos: al asignar a cada hilo una tabla
privada en lugar de una tabla compartida con \emph{mutex}, la fase paralela quedó
libre de contención.

\subsection{¿Agregar más hilos garantiza siempre un mejor rendimiento?}

\textbf{No.} Los datos obtenidos lo evidencian en tres niveles:

\paragraph{(a) El rendimiento mejora, pero cada vez menos.}
Al duplicar de 2 a 4 hilos en el archivo pequeño, el tiempo descendió de
\num{0.002188}~s a \num{0.001745}~s: \textbf{se duplicaron los recursos para
obtener apenas un 20\,\% de mejora}. El segundo par de hilos aportó una fracción
mínima de lo que aportó el primero.

\paragraph{(b) La eficiencia se degrada severamente.}
Pasar de 2 a 4 hilos redujo la eficiencia del 81.76\,\% al 51.26\,\% en el archivo
pequeño. En términos prácticos, con 4 hilos \textbf{casi la mitad de la capacidad
de cómputo asignada se desperdició} en coordinación.

\paragraph{(c) Existen ejecuciones donde 4 hilos resultó más lento que 2.}
No se trata únicamente de un promedio que se degrada. En el archivo pequeño, la
corrida~5 registró \num{0.001376}~s con 2 hilos frente a \num{0.001426}~s con 4, y
la corrida~4 fue prácticamente un empate (\num{0.001510}~s frente a
\num{0.001490}~s). Agregar hilos en esos casos \textbf{no aportó absolutamente
nada}.

\paragraph{Conclusión.}
El número óptimo de hilos depende de la relación entre trabajo útil y
\emph{overhead}, no de la cantidad de núcleos disponibles. El procesador empleado
dispone de 10 núcleos, pero con \texttt{archivo.txt} la eficiencia ya se situaba
por debajo del 52\,\% con 4 hilos; emplear 10 hilos sobre 24~KB de texto sería
contraproducente, dado que el costo de crear y fusionar diez tablas superaría el
ahorro obtenido en el conteo.

% ====================================================== 7. CONCLUSIONES
\section{Conclusiones}

\begin{enumerate}[label=\textbf{\arabic*.}]
  \item La paralelización del conteo de frecuencia \textbf{sí produce una mejora
        real y medible}: hasta $2.89\times$ con 4 hilos, manteniendo resultados
        idénticos a la versión secuencial en todas las configuraciones probadas.
  \item El \textbf{mayor \emph{speedup}} se obtuvo con 4 hilos, pero la
        \textbf{mayor \emph{efficiency}} con 2 hilos. Ambas métricas responden a
        criterios de optimización distintos y no deben confundirse.
  \item El \emph{speedup} crece de forma \textbf{sublineal} y la \emph{efficiency}
        \textbf{decrece} al agregar hilos, comportamiento explicado por la ley de
        Amdahl: la reducción de las tablas es estrictamente secuencial y su costo
        aumenta con el número de hilos.
  \item El \textbf{tamaño del problema resultó ser el factor más influyente}. El
        mismo código con 4 hilos pasó del 51\,\% al 72\,\% de eficiencia
        únicamente al incrementar el tamaño del archivo. El beneficio del
        paralelismo depende de que exista trabajo suficiente para amortizar el
        costo de coordinación.
  \item La decisión de emplear \textbf{tablas privadas por hilo} en lugar de una
        tabla compartida con \emph{mutex} fue determinante. Concentra todo el
        costo de sincronización en un único punto —la reducción— en lugar de
        distribuirlo a lo largo de toda la ejecución.
\end{enumerate}

% ====================================================== 8. REPRODUCCION
\section{Reproducción de los resultados}

\begin{lstlisting}[style=sh,caption={Comandos para reproducir el experimento}]
make                # compila ambas versiones
make verificar      # verifica que la salida paralela sea identica
make benchmark      # 5 corridas por configuracion, speedup y efficiency

# Compilacion manual:
gcc -O2 -Wall          -o conteo_secuencial conteo_frecuencia_secuencial.c
gcc -O2 -Wall -pthread -o conteo_paralelo   conteo_frecuencia_paralelo.c

# Ejecucion individual:
./conteo_secuencial archivo.txt
./conteo_paralelo   archivo.txt 4
./conteo_paralelo   archivo.txt 4 --silencioso   # solo el tiempo

# Experimento complementario:
make benchmark ARCHIVO=archivo_grande.txt REPS=5
\end{lstlisting}

\noindent
Los tiempos crudos de cada corrida se encuentran en los archivos
\texttt{resultados\_archivo.csv} y \texttt{resultados\_grande.csv}.

\vfill
\begin{center}
\color{grisln}\rule{0.5\linewidth}{0.6pt}
\end{center}

\end{document}
